\chapter{Introduction}\label{ch:introduction}

The inverse problem requires recovering physical parameters
from indirect measurements. This problem is ubiquitous throughout the
sciences; occurring whenever the quantities of interest cannot be observed
directly. In electrochemistry, the parameters governing a reaction at
an electrode surface must be inferred from the current response to an
applied potential. This inference is the electrochemical inverse
problem.

Predicting the current response for a given set of parameters is the
corresponding forward problem, which in the present context, reduces to
solving a diffusion PDE with kinetic boundary conditions
\cite{compton2014understanding}. Solving the inverse problem then
amounts to finding the parameters whose forward solution best
explains the observed data. The difficulty is threefold: the
mapping from parameters to current is nonlinear, distinct parameter
configurations can produce similar voltammograms, and in certain
kinetic regimes some parameters cease to influence the current altogether.
Point estimates cannot capture these difficulties. They provide no
account of uncertainty, cannot distinguish a well-constrained parameter
from one that is insensitive to the data, and do not reveal correlations
between parameters. Bayesian inference addresses this by computing the
posterior distribution conditioned on the observed
data \cite{gavaghan2018bayesian}.

The standard tool for sampling such posteriors is Markov chain Monte
Carlo. In this setting each likelihood evaluation requires a PDE solve,
so every proposed sample, whether accepted or rejected, is expensive.
Gradient-free samplers such as random-walk Metropolis--Hastings (RWMH)
explore the parameter space by undirected local proposals and
consequently waste many of these solves on rejected candidates,
particularly as the number of parameters grows. Gradient information
helps in two ways: it can initialise the sampler near a posterior mode
via optimisation; and it can inform the sampling itself
through Hamiltonian Monte Carlo (HMC), which uses the gradient to
construct long-range, low-rejection proposals.

The challenge is obtaining these gradients cheaply. Finite differences
require at least one additional forward solve per parameter and introduce
truncation error. Whereas, reverse-mode automatic differentiation computes exact
gradients at a cost bounded by a small multiple of the forward solve,
independent of parameter dimension \cite{griewank2008evaluating}.
Na\"ively taping an entire PDE solver, however, records every
intermediate of the linear solve -- elimination multipliers, modified
diagonals, partial sums -- at each timestep. We instead derive a
custom adjoint rule for the banded linear solve and
register it with the automatic differentiation framework, \emph{JAX}
\cite{jax2018github}, so that gradient computation uses only the matrix
diagonals and solution vector.

This work sits within a broader shift toward differentiable scientific
computing; Kitchin et~al.\ \cite{kitchin2025beyond} characterise
differentiable programming as a fifth paradigm for modelling in chemical
engineering. Concurrently, Chen et~al.\ \cite{chen2026differentiable}
have developed differentiable electrochemical simulators for
gradient-based parameter recovery. Our work is complementary: we
adopt a Bayesian framework rather than optimisation alone, and we
compute gradients via a custom adjoint rather than taping the full
solver. The Bayesian approach to voltammetric inference via MCMC was
established by Gavaghan et~al.\ \cite{gavaghan2018bayesian}; we
extend their framework by replacing the gradient-free sampler with
gradient-informed alternatives enabled by the differentiable solver.

This dissertation makes three contributions. First, we develop a
differentiable forward solver for voltammetric PDEs, in which a custom
adjoint rule for the banded linear system provides exact gradients at low
memory cost and generalises to the wider banded systems arising in
multi-species reactions. Second, we integrate this solver with the
No-U-Turn sampler (NUTS)~\cite{hoffman2014no}, an
extension of HMC, to obtain posterior distributions over the physical
parameters. Third, we apply the framework to three reaction
mechanisms of increasing complexity:  a single electron transfer
(3~parameters), a heterogeneous reaction (7~parameters), and an
adsorption reaction (9~parameters). We find that the advantage of
NUTS over RWMH grows with the dimensionality of the parameter space.
The methodology is developed through the single-electron-transfer
reaction, which serves as a running example; the final two chapters
apply the complete framework to the harder reactions concisely.
